\section{The complete cylinder realization}\label{sec:operator}

\subsection{Hilbert space, maximal domain and spectrum}
Let \(\Sigma\) be the spinor bundle on the unit sphere and set
\[
 \KK=L^2(S^2,\Sigma;\dd\Omega),\qquad
 \HH_C=L^2(\R_v;\KK),\qquad
 g_C=\dd v^2+\tfrac14g_{S^2}.
\]
We use product measure \(\dd v\dd\Omega\). Metric-volume fields have
measure \(\dd v\dd\Omega/4\); multiplication by two is the unitary
conversion, and the readout must then be composed with its inverse.

The unit-sphere Dirac operator \(D\) has eigenvalues \(\pm n\),
\(n\geq1\), with multiplicity \(2n\) for each sign
\cite{CamporesiHiguchi1996}. Put \(A=D^2\) and let \(P_n\) be its
complete shell projection, of rank \(4n\). The cylinder square and its
relative-energy shift are
\begin{equation}\label{eq:K}
 L_0=-\partial_v^2+4A,\qquad
 K=L_0-4I=-\partial_v^2+4(A-I).
\end{equation}
Use the unitary Fourier convention
\(\widehat f(p)=(2\pi)^{-1/2}\int\ee^{-\ii pv}f(v)\dd v\).
The complete Fourier/shell decomposition identifies \(K\) with the
real multipliers
\[
 m_n(p)=p^2+4(n^2-1),\qquad
 \HH_C\simeq\bigoplus_{n\geq1}L^2(\R_p;\Ran P_n).
\]

\begin{theorem}[Maximal realization]\label{thm:operator}
The operator
\begin{equation}\label{eq:domain}
 \Dom K=\left\{\Psi\in\HH_C:
 \sum_{n\geq1}\int_\R m_n(p)^2\norm{P_n\widehat\Psi(p)}_\KK^2\dd p<\infty\right\},
 \quad \widehat{K\Psi}_n=m_n\widehat\Psi_n,
\end{equation}
is nonnegative and self-adjoint. Its domain equals \(\Dom L_0\), with
equivalent graph norms. A common core is
\begin{equation}\label{eq:core}
 \mathcal C=\bigcup_{N\geq1}C_c^\infty
 \left(\R;\bigoplus_{n=1}^N\Ran P_n\right).
\end{equation}
The spectrum of \(K\) is purely absolutely continuous and equals
\([0,\infty)\); zero is not an eigenvalue.
\end{theorem}
\begin{proof}
Finite shells with bounded momentum support give a dense domain.
Reality of the multipliers proves symmetry and positivity. If
\(\eta\in\Dom K^*\) has adjoint image \(\zeta\), test against arbitrary
\(L^2\) fields supported in finitely many shells and a bounded momentum
interval. There the multipliers are bounded, so the adjoint identity
forces \(\zeta_n=m_n\eta_n\) almost everywhere. Exhausting the shells
and intervals shows precisely the integrability in \eqref{eq:domain}.
The converse follows by Cauchy--Schwarz. Hence \(K^*=K\).

The multiplier of \(L_0\) is \(m_n+4\), and
\(m_n^2\leq(m_n+4)^2\leq2m_n^2+32\). This proves domain equality
and graph-norm equivalence. Truncating to complete shells converges in
the graph norm. On finitely many shells the graph norm is equivalent
to the radial \(H^2\) norm. Multiplication by a smooth cutoff
\(\eta(v/R)\), followed by mollification, approximates each component
in \(H^2\) by a compactly supported smooth function. Indeed the cutoff
derivative terms have factors \(R^{-1}\) and \(R^{-2}\), and all
remaining tails vanish in \(L^2\). This proves \eqref{eq:core}.

Outside \([0,\infty)\), the reciprocal multipliers give a bounded
resolvent. The first shell provides approximate spectral vectors near
every value \(p^2\geq0\), so the spectrum is the stated interval.
Each fixed-energy level has measure zero in every shell, ruling out
eigenvalues. On each open momentum half-line, the change of variable
\(p\mapsto p^2+4(n^2-1)\) gives an absolutely continuous spectral
measure. The thresholds have no atoms; the countable orthogonal sum
remains absolutely continuous.
\end{proof}

Both Dirac signs are retained. In a signed cylinder block, the first-order
operator has matrix \(\bigl(\begin{smallmatrix}2n&p\\p&-2n\end{smallmatrix}\bigr)\),
with eigenvalues \(\pm\sqrt{p^2+4n^2}\), each of angular multiplicity
\(2n\). Thus \(L_0\) has spectrum \([4,\infty)\), and subtracting its
threshold does not select a sign or create a normalizable ground state.

\subsection{Heat evolution and its domain}
Define \(S(\tau)\), \(\tau\geq0\), by multiplication with
\(\ee^{-\tau m_n(p)}\) in every shell.

\begin{proposition}\label{prop:semigroup}
The family \(S(\tau)\) is a strongly continuous contraction semigroup
with norm one and generator \(-K\) on exactly \(\Dom K\). For every
integer \(j\geq1\),
\begin{equation}\label{eq:smoothing}
 \norm{K^jS(\tau)}\leq\left(\frac{j}{\ee\tau}\right)^j\quad(\tau>0).
\end{equation}
No bounded reverse-time semigroup exists on \(\HH_C\).
\end{proposition}
\begin{proof}
The semigroup law and contraction bound are multiplier identities.
The norm is one because the first-shell momentum can approach zero.
Dominated convergence proves strong continuity. On \(\Dom K\),
\(|(\ee^{-hm}-1)/h|\leq m\) for \(h>0\), so the strong right
derivative is \(-K\). Conversely, boundedness of the squared norm of
the difference quotients and Fatou's lemma force the domain integral
in \eqref{eq:domain}. The supremum of \(m^j\ee^{-\tau m}\) is
\((j/(\ee\tau))^j\). Finally \(\ee^{\delta p^2}\) is unbounded
already in the first shell for every \(\delta>0\).
\end{proof}

The domain of \(K^j\) is given by the same decomposition with weight
\(m_n^{2j}\), together with underlying \(L^2\) membership.
The core argument extends using \(H^{2j}\). For the complete source,
every prepared power has the explicit bound:
\begin{equation}\label{eq:prepared-power}
 \norm{K^jE_\chi F_\lambda}^2
 =\int_\R |u|^{4j}|F_\lambda(u)|^2\dd u
 \leq B_0^2J_{4j,4}(2L,1).
\end{equation}
Here the encoder is defined next; the identity uses its first-shell
support. These source-specific bounds do not supply a bounded inverse
flow on arbitrary ambient fields.

\subsection{Preparation and absolute observation}
Fix a unit \(\chi\in\Ran P_1\), allowing both Dirac signs, and set
\begin{equation}\label{eq:encoder}
 \widehat{E_\chi f}(p)=f(p)\chi,\qquad
 J_\chi\Psi(x)=\inner{\chi}{\Psi(x)}_\KK,\qquad
 R_\chi=\sqrt{\pi/2}\,J_\chi.
\end{equation}
The source variable is calibrated to momentum \(p=u\); the observed
variable is position \(x=v\). Pairings are conjugate-linear in their
first argument, so the readout is linear in the field.

\begin{theorem}[Exact calibrated realization]\label{thm:realization}
The encoder is isometric and \(R_\chi:\HH_C\to L^2(\R_x)\) is bounded
with norm \(\sqrt{\pi/2}\). If \(M\) denotes multiplication by \(u^2\)
and \(B=-\partial_x^2\) on \(H^2(\R)\), then
\begin{align}
 KE_\chi f&=E_\chi Mf,\quad
 E_\chi f\in\Dom K\ \Longleftrightarrow\ u^2f\in L^2,
 \label{eq:intertwine-domain}\\
 R_\chi K\Psi&=BR_\chi\Psi\quad(\Psi\in\Dom K),\label{eq:readout-domain}\\
 S(\tau)E_\chi&=E_\chi\ee^{-\tau M},\qquad
 R_\chi S(\tau)=\ee^{-\tau B}R_\chi.\label{eq:intertwine-heat}
\end{align}
In particular
\begin{equation}\label{eq:realization}
 R_\chi S(\tau)E_\chi F_{\lambda_0}=H_{\lambda_0-\tau},\qquad
 \Psi_\tau:=S(\tau)E_\chi F_{\lambda_0}
 =\sqrt{2/\pi}\,h(\tau,\cdot)\chi.
\end{equation}
\end{theorem}
\begin{proof}
Plancherel and \(\norm\chi=1\) prove isometry. Angular Cauchy--Schwarz
gives \(\norm{J_\chi}\leq1\), with equality on fields \(f(x)\chi\).
The first-shell multiplier is \(p^2\), proving
\eqref{eq:intertwine-domain}. For \(\Psi\in\Dom K\),
\[
 \int p^4|\inner\chi{\widehat\Psi(p)}|^2\dd p
 \leq\sum_n\int m_n(p)^2\norm{P_n\widehat\Psi(p)}^2\dd p.
\]
Thus \(R_\chi\Psi\in H^2\), and pairing the multiplier with
\(\chi\) proves \eqref{eq:readout-domain}. The bounded semigroup
identities follow directly without cancelling unbounded operators.
For \(f\in L^1\cap L^2\), Fourier inversion gives
\[
 R_\chi E_\chi f(x)=\frac{\sqrt{\pi/2}}{\sqrt{2\pi}}
 \int_\R f(u)\ee^{\ii xu}\dd u
 =\frac12\int_\R f(u)\ee^{\ii xu}\dd u.
\]
Apply this to the complete even source. The smooth representative and
all fixed derivatives in \eqref{eq:realization} follow from
Proposition~\ref{prop:source-regular} and \eqref{eq:prepared-power}.
\end{proof}

The rank-one angular preparation is a reducing channel of the full
operator, not an alteration of its ambient spectrum. The map \(R_\chi\)
is bounded into \(L^2\); point evaluation on an arbitrary field is not
claimed bounded. The supplied source moments select the pointwise
representative used in the collision analysis.

\subsection{The energy origin and the exact scalar baseline}
The absolute calibration makes the energy shift visible:
\[
 R_\chi\ee^{-\tau L_0}E_\chi F_{\lambda_0}
 =\ee^{-4\tau}H_{\lambda_0-\tau}.
\]
It cannot be removed within the following fixed bounded-readout class.

\begin{proposition}[Scoped unshifted obstruction]\label{prop:unshifted}
No fixed \(\Psi_0\in\HH_C\) and fixed bounded linear
\(R:\HH_C\to L^2(\R)\) satisfy
\(R\ee^{-\tau L_0}\Psi_0=H_{\lambda_0-\tau}\) for every \(\tau\geq0\).
\end{proposition}
\begin{proof}
The left-hand norm is at most \(\norm R\norm{\Psi_0}\ee^{-4\tau}\).
On the other hand Plancherel and \(u=y/\sqrt\tau\) give
\[
 \tau^{1/2}\norm{H_{\lambda_0-\tau}}_2^2
 =\frac\pi2\int_\R\ee^{-2y^2}
 |F_{\lambda_0}(y/\sqrt\tau)|^2\dd y
 \longrightarrow\Phi(0)^2(\pi/2)^{3/2}.
\]
Boundedness and continuity of the source justify dominated convergence.
Hence the target norm is asymptotic to
\(\Phi(0)(\pi/2)^{3/4}\tau^{-1/4}\), with positive coefficient,
contradicting exponential decay.
\end{proof}
The proposition concerns the entire half-line and the stated bounded
linear class. It does not exclude finite-interval realizations,
unbounded observations or time-dependent calibration.

Finally, for arbitrary \(f,g\in L^2(\R)\),
\begin{equation}\label{eq:exact-error}
 \norm{R_\chi S(\tau)E_\chi(f-g)}_2
 =\sqrt{\pi/2}\,\norm{\ee^{-\tau u^2}(f-g)}_2.
\end{equation}
This is exactly the direct scalar Fourier error identity. The operator
realization supplies a well-defined field interpretation and domain
theory, while this comparison establishes no estimation gain from the
representation alone.
